%%=============================================================================
%% Methodologie
%%=============================================================================

\chapter{\IfLanguageName{dutch}{Methodologie}{Methodology}}%
\label{ch:methodologie}

%% TODO: In dit hoofstuk geef je een korte toelichting over hoe je te werk bent
%% gegaan. Verdeel je onderzoek in grote fasen, en licht in elke fase toe wat
%% de doelstelling was, welke deliverables daar uit gekomen zijn, en welke
%% onderzoeksmethoden je daarbij toegepast hebt. Verantwoord waarom je
%% op deze manier te werk gegaan bent.
%% 
%% Voorbeelden van zulke fasen zijn: literatuurstudie, opstellen van een
%% requirements-analyse, opstellen long-list (bij vergelijkende studie),
%% selectie van geschikte tools (bij vergelijkende studie, "short-list"),
%% opzetten testopstelling/PoC, uitvoeren testen en verzamelen
%% van resultaten, analyse van resultaten, ...
%%
%% !!!!! LET OP !!!!!
%%
%% Het is uitdrukkelijk NIET de bedoeling dat je het grootste deel van de corpus
%% van je bachelorproef in dit hoofstuk verwerkt! Dit hoofdstuk is eerder een
%% kort overzicht van je plan van aanpak.
%%
%% Maak voor elke fase (behalve het literatuuronderzoek) een NIEUW HOOFDSTUK aan
%% en geef het een gepaste titel.

Dit onderzoek bestaat uit een analyse van de bedrijfscasus en een vergelijkende studie met een proof of concept.
Eerst worden de huidige werkwijze en de bedrijfsvereisten onderzocht. Daarna worden twee implementaties
uitgewerkt en onder dezelfde omstandigheden getest. De resultaten vormen de basis voor een advies aan het bedrijf.

In de eerste fase wordt gekeken hoe werknotities en materiaalregistraties momenteel met Installware worden verwerkt.
Om dit te achterhalen wordt contact opgenomen met co-promotor Jarne Schagen. De vragen gaan over de gegevens die medewerkers
registreren, de mogelijkheden tijdens netwerkuitval, de stappen
na verbindingsherstel en situaties waarin meerdere medewerkers dezelfde werkbon aanpassen. 
Ter voorbereiding van de bevraging is een vragenlijst opgesteld voor co-promotor Jarne Schagen. De vragen zijn in niet-technische taal
geformuleerd en vragen naar concrete voorbeelden uit de dagelijkse werking. Ze behandelen het gebruik van werkbonnen, werknotities en
materiaalregistraties, de mogelijkheden bij netwerkuitval en het aanpassen van dezelfde werkbon door meerdere medewerkers.
De antwoorden dienen als basis om de huidige werkwijze en bedrijfsvereisten verder in kaart te brengen.

Het proof of concept bestaat uit een eenvoudige mobiele testapplicatie met lokale opslag en een backend.
Als voorlopig technisch uitgangspunt wordt React Native met Expo en
TypeScript gebruikt~\autocite{ReactNativeTypeScript}, met SQLite voor lokale opslag~\autocite{ExpoSQLite}. Voor de
CRDT-gebaseerde implementatie wordt Yjs gebruikt~\autocite{YjsDocumentation}. De eerste implementatie bewaart wijzigingen in een lokale
wachtrij. De backend controleert via versienummers of een wijziging op verouderde gegevens is gebaseerd. Conflicten worden
zichtbaar gemaakt en volgens de vooraf vastgelegde regels afgehandeld. De tweede implementatie gebruikt CRDT's om wijzigingen
samen te voegen. Beide varianten ondersteunen dezelfde handelingen voor werknotities en materiaalregistraties. Ze gebruiken
zoveel mogelijk dezelfde gebruikersinterface, gegevens en testomgeving, zodat er zo weinig mogelijk verschillen zijn die invloed zouden kunnen hebben.

De tests gebruiken fictieve werkbonnen die aansluiten bij de vastgestelde bedrijfsvereisten. Er worden scenario's opgesteld voor
herstel van de verbinding, een nieuwe onderbreking tijdens synchronisatie, herhaalde verzending en gelijktijdige
wijzigingen door meerdere gebruikers.

Gegevenscorrectheid wordt beoordeeld door de gegevens te vergelijken met de vooraf vastgelegde verwachte resultaten.
Daarbij worden ontbrekende registraties, ongewenste duplicaten, afwijkingen van de werkbonregels en verschillen tussen gegevenskopieën
bijgehouden. De synchronisatietijd wordt gemeten
vanaf het herstel van de verbinding tot alle testwijzigingen zijn verwerkt. Elk testgeval wordt dertig keer uitgevoerd vanuit dezelfde
beginsituatie. De resultaten worden samengevat met de mediaan en de spreiding.

Tot slot worden beide implementaties naast de bedrijfsvereisten gelegd. Het advies vermeldt zowel de voordelen als de beperkingen
van elke aanpak.


